% Français 9115, batch 15 — Gallica views 281–300.
% Pass: Opus 5 (claude-opus-5), 2026-09-19 — first pass, unchecked against
% the views by a human.
%
% What the twenty views hold. The Latin fair copy of Euler's « Investigatio
% motuum quibus laminæ et virgæ elasticæ contremiscunt », which filled
% batches 8 to 10, is over. These twenty views are a memoir of her own, in
% French, in her regular fair hand but heavily corrected between the lines:
% a critical examination of the very paragraph of Euler she had just copied
% — « N.o 47 », the rod pinned at an intermediate point, copied in Latin at
% views 192 and following. Ten views carry text and ten are blank: the
% leaves are written on the recto only, and every even view of the batch
% (284, 286, 288, 290, 292, 294, 296, 298, 300) is the blank verso of the
% recto before it, showing nothing but the recto's writing through the
% paper, reversed. There are no repeats of a photographed page, and no
% microfilm target.
%
% The text runs on without a break from one written leaf to the next, and
% it is numbered by paragraph, 19 to 23, the figure written inside the
% line. It does continue a sentence begun before this batch: view 283 opens
% on « de sorte que pour une valeur quelconque de x on peut toujours… ».
%
%   v. 281      not a new leaf: the leaf of view 280 photographed a second
%               time with its large béquet turned up, which uncovers the six
%               lines at the foot of the page that the béquet hid. The
%               passage is struck through with a large cross.
%   v. 283      end of the demonstration that, the two ends being simply
%               supported, a node at the distance ax from one end forces a
%               node at ax from the other, the equations being homogeneous
%               functions of ax and a−ax; then § 19, her claim that all of
%               this can be reached « sans avoir recours a la méthode dont
%               il s'agit ici », by the ordinate alone set equal to zero at
%               a node.
%   v. 285–287  the same conclusion carried through for the free ends and
%               for the fixed ends, by adding and subtracting the equation
%               of the half-length; § 20 (v. 287), the reproach itself: the
%               few conclusions Euler's four equations yield are got more
%               simply from the ordinate, and his equations « semblent
%               indiquer des conclusions qui ne sont pourtant pas
%               admissibles ». She states what he assumes at N.o 47 — that
%               the pin interrupts the continuity of the curve, the two
%               portions being given by equations differing in their
%               coefficients.
%   v. 289      the argument against it: the motion with the pin is still
%               one of the regular motions of the rod, hence one of those
%               the six cases determine, hence continuous; therefore
%               α'=α, β'=β, γ'=γ, δ'=δ. Then, in so many words, « Je n'ose
%               admettre ce resultat contraire a l'opinion d'Euler […] sans
%               le soumettre a un examen plus approfondi », and the
%               announcement of the proposition she will prove. The body of
%               the leaf is crossed out and rewritten in the left margin.
%   v. 291      § 21, the proposition itself, and the proof that one
%               equality of coefficients entails the others; the case of two
%               free ends reduced to a single identity to be verified.
%   v. 293–295  that identity developed and verified, first for sin ω
%               positive, then for sin ω negative.
%   v. 297      § 22, the same for two fixed ends: the two expressions of
%               α/α' differ only by the sign of e^{−xω}, and the change of
%               sign is compensated.
%   v. 299      § 23, the case of two supported ends, where Euler's
%               α/α' = (e^{xω} − e^{(2−x)ω})/(e^{xω} − e^{−xω}) settles
%               nothing because every coefficient but δ vanishes and the
%               ratio is indeterminate; the eight coefficients of Euler's
%               § 50 quoted (they are the ones she had copied in Latin at
%               view 200), and the conclusion that « l'appareille des quatre
%               équations trouvées par Euler est au moins inutile ». The
%               leaf stops mid-sentence on « les deux »; the sequel is not
%               in this batch.
%
% The numbers on the leaves, and why no \folio{} is written. Two series,
% one per written leaf, both apparently in ink, and none on the blank
% versos. (1) At the head, in the middle, underlined, in her own hand and
% her own ink: 30 at view 283, then 31, 32, 33, 34, 35, 36, 37, 38 at views
% 285 to 299, without a gap. (2) At the top right of the same leaves, in a
% larger and rounder hand and a broader pen: 138 at view 283, then 139 to
% 146 at views 285 to 299, again without a gap, one to a leaf. Its place —
% top right of a recto, one figure per leaf, one series — is the place of
% the Bibliothèque's foliation, but the figures are penned, not the thin
% grey pencil the foliation usually is, and they are of a piece with the
% ink series recorded in batches 1 to 10 of this volume (89 to 97 at views
% 180 to 199). No pencilled foliation was read anywhere in the batch, so no
% \folio{} is written here, as in every batch of this volume so far. Every
% number above was read on its view; none is inferred. View 281 carries no
% legible number.
%
% Notation. Hers, kept: $x$ is a fraction of the length, so that the point
% of support is at the distance $ax$ from one end and $a-ax$ from the
% other, and $\omega$ is the angle of the six cases; $\alpha$, $\beta$,
% $\gamma$, $\delta$ and $\alpha'$, $\beta'$, $\gamma'$, $\delta'$ are
% Euler's eight coefficients, one set for each portion of the rod. Her
% $\gamma$ is written like a « v », as a pass already settled on the Latin
% copy. « sin », « cos » and « tang. » are written with a long s and are
% set here as \text{sin}, \text{cos}, \text{tang.}; the stop after them is
% irregular on the leaf and has not been recorded stop by stop. Braces
% $\{\ \}$ are hers and are kept where she wrote them. Spelling is hers and
% is not flagged: « coëfficiens », « extrimités », « parconséquent »,
% « parséquent », « dela », « demême », « adire », « apresent »,
% « inconcilliables », « etres », « l'appareille », « aiégard »,
% « corrispondant », « atribue », « fesant », « suivant » for « suivent ».
% Where the leaf and the calculation disagree the leaf stands, with a note:
% view 285 (« $x-1$ » for $1-x$), view 291 (an exponent without its sign,
% another without its $\omega$), view 295 (the same exponent twice), view
% 297 (« les équations » for « les extremités », and
% « $e^{(1-x)\omega}$ » for $e^{(x-1)\omega}$).
%
% Her way of correcting. The additions are written above or below the line,
% almost never with a caret, and they often replace a word that is struck
% on the line: they are given here as \add{} beside the \struck{} they
% replace, in the order the eye takes them. Long horizontal strokes run
% through many words that make perfect sense; following what an earlier
% pass settled on this hand, those are her joining strokes and are not
% given as deletions — only a separate, heavier stroke is. Two whole pages
% are crossed with a large pen cross (views 281 and 289) and are said to be
% so in a note rather than wrapped in \struck{}. The marginal notes are
% hers; they are placed against the line they face, with the sign (an
% asterisk, a cross) that calls them where she wrote one.
%
% What could not be read. \ill{} stands in ten places: three struck groups
% at view 283 (one after « on l'a vu », whose last word ends in « -ait »,
% and two on the line « l'état des extrémités », the second followed by
% « de »); a blotted sign between two products at view 285; two short
% struck groups at view 291; a blotted sign at view 293; at view 295 a
% struck and repassed group and a blotted sign; a short struck word after
% « au lieu de » at view 297. The character struck after « que » at view
% 281 is described in its note rather than marked, being caught in the
% cross that cancels the page. Readings offered with doubt are marked
% \uncertain{}: « cesse
% d'être » at view 285; « quinzième » at view 287; « demême » at view 295;
% « ici en question » at view 297; « par » at view 299. Nothing was guessed
% to fill a gap.
%
% Nothing here was checked against any published edition of Euler's memoir
% or of Sophie Germain's papers.
\documentclass[11pt,a4paper]{article}
% The preamble shared by every transcription of the mathematicians' manuscripts in Gallica.
%
% It rests on one idea: **uncertainty compounds, it does not resolve.** A
% transcription that smooths over a doubtful word has destroyed the only thing
% it was made for — a reader must be able to tell, at every point, what was on
% the page and what was supplied. The apparatus macros below are therefore the
% critical apparatus, not decoration.
%
% The same commands are recognised by `scripts/render.mjs`, which produces the
% site's reading view, and by `scripts/tei.mjs`. A macro added here must be
% added there too, or rendering fails loudly — which is the wanted behaviour.
%
% Comments here are English, like the rest of the repository. The documents
% this preamble sets are French, and that is deliberate: see the skills.

\ProvidesPackage{germain}[2026/09/15 Transcriptions of mathematicians' manuscripts in Gallica]

\usepackage{amsmath,amssymb,amsthm}
\usepackage{mathrsfs}
% Kept from the parent preamble so the renderer's diagram support stays
% available; these manuscripts draw no commutative diagrams, and nothing here uses it.
\usepackage{tikz-cd}
\usepackage[english,french]{babel}
\usepackage{fontspec}

% --- Greek ------------------------------------------------------------------
% The text face stays fontspec's default, Latin Modern, which has no Greek:
% XeTeX drops every glyph it lacks with a line in the log and nothing on the
% page, and the Greek words the manuscripts quote vanished from the PDFs. So
% the Greek and Greek Extended blocks get a XeTeX character class of their own,
% and entering or leaving a run of it switches the family to CMU Serif — the
% same Computer Modern design, with polytonic Greek — and back. Only the
% family changes: a Greek word in \textit or \textbf keeps its shape and series.
% The transitions to and from every other class carry the tokens that class
% has towards an ordinary letter, so French spacing before « ; » or inside
% « » still holds after a Greek word. No grouping, so no brace or \endgroup
% can come out unbalanced; maths is left alone, since XeTeX inserts nothing
% there. Kept identical in `exercices.sty`.
\makeatletter
\newfontfamily\germainGreekfont{cmun}[
  Extension=.otf, NFSSFamily=germaingreek,
  UprightFont=*rm, ItalicFont=*ti, SlantedFont=*sl,
  BoldFont=*bx, BoldItalicFont=*bi, BoldSlantedFont=*bl]
\newXeTeXintercharclass\germain@greek
\def\germain@greekrange#1#2{%
  \@tempcnta=#1\relax
  \loop
    \XeTeXcharclass\@tempcnta=\germain@greek
    \ifnum\@tempcnta<#2\advance\@tempcnta\@ne\repeat}
\germain@greekrange{"0370}{"03FF}
\germain@greekrange{"1F00}{"1FFF}
\def\germain@greekprev{\rmdefault}
\def\germain@greekin{%
  \edef\germain@greekprev{\f@family}\fontfamily{germaingreek}\selectfont}
\def\germain@greekout{\fontfamily{\germain@greekprev}\selectfont}
\def\germain@greekjoin{%
  \XeTeXinterchartokenstate=\@ne
  \@tempcnta=\z@
  \loop
    \ifnum\@tempcnta=\germain@greek\else
      \XeTeXinterchartoks\germain@greek\@tempcnta=\expandafter{%
        \expandafter\germain@greekout\the\XeTeXinterchartoks\z@\@tempcnta}%
      \XeTeXinterchartoks\@tempcnta\germain@greek=\expandafter{%
        \the\XeTeXinterchartoks\@tempcnta\z@\germain@greekin}%
    \fi
    \ifnum\@tempcnta<4095\advance\@tempcnta\@ne\repeat}
% babel-french sets its own classes, and saves and restores the interchar
% state, each time a language is selected: join again after it does.
\addto\extrasfrench{\germain@greekjoin}
\addto\extrasenglish{\germain@greekjoin}
\AtBeginDocument{\germain@greekjoin}
\makeatother

\usepackage{geometry}
\geometry{a4paper,margin=25mm,marginparwidth=28mm}
\usepackage{marginnote}
\usepackage{xcolor}
\usepackage[normalem]{ulem}
\setlength{\emergencystretch}{3em}
\usepackage{hyperref}
\hypersetup{colorlinks=true,linkcolor=black,urlcolor=black,pdfborder={0 0 0}}

% --- Batch metadata ---------------------------------------------------------
% Taken as they stand by the rendering script, which shows them at the head of
% the reading view. They fix what the file claims to cover. The macro names are
% the parent project's, so that the scripts stay shared; \folder here means the
% volume's slug — `fr-9115` — and \pages counts Gallica views.

\newcommand{\folder}[1]{\gdef\grothFolder{#1}}
\newcommand{\batch}[1]{\gdef\grothBatch{#1}}
\newcommand{\pages}[2]{\gdef\grothFirst{#1}\gdef\grothLast{#2}}
\newcommand{\foldertitle}[1]{\gdef\grothFoldertitle{#1}}

% The holder's own shelfmark, printed as they write it: « Français 9115 ».
\newcommand{\shelfmark}[1]{\gdef\grothShelfmark{#1}}

% The dating, copied verbatim from the holder's catalogue — never inferred
% here. « XIXe siècle » is the BnF's; a pass that dates a leaf from its content
% says so in a \note{}, as its own inference.
\newcommand{\dating}[1]{\gdef\grothDating{#1}}

% The status notice, printed in the title block and repeated as a diagonal
% watermark on every page of the PDF. The manuscripts are in the public
% domain; what this line declares is not a right but a quality — an
% unverified first pass by a machine. The skills write the line; a file
% without it gets no watermark, so leaving it out is visible in the output.
\newcommand{\watermark}[1]{\gdef\grothWatermark{#1}}

\gdef\grothFolder{?}\gdef\grothBatch{}\gdef\grothFirst{?}\gdef\grothLast{?}%
\gdef\grothFoldertitle{}\gdef\grothShelfmark{}\gdef\grothDating{}\gdef\grothWatermark{}

% --- The résumé -------------------------------------------------------------
\newenvironment{resume}
  {\small\setlength{\parskip}{0.45em}}
  {\par\medskip\hrule\medskip}

% --- Keywords ---------------------------------------------------------------
% In English, closing the modernised reading's résumé: the modern vocabulary
% under which to search for what the leaves construct. The manifest extracts
% them to tag the volume — this line is the single source of the tags.
\newcommand{\keywords}[1]{\par\smallskip\noindent
  {\footnotesize\textsc{Keywords} --- \textit{#1}}\par}

% --- The critical apparatus -------------------------------------------------

% The Gallica view number, in the margin. This is the anchor that lets a
% disagreement over a formula be settled by looking at *one* image rather than
% a volume of 750. The site uses it to turn the facsimile as the transcription
% is scrolled. A view is one scanned image: usually one side of a leaf.
\newcommand{\page}[1]{%
  \marginnote{\footnotesize\color{gray!70}\textsf{v.\,#1}}%
  \label{page:#1}\ignorespaces}

% The BnF's pencilled foliation, where the leaf carries it and a pass can read
% it: « 198r ». This is what the literature cites — « ff. 198r–208v » — and
% the one line that turns a citation into a view. Recorded, never inferred:
% a folio a pass cannot read is not written.
\newcommand{\folio}[1]{%
  \marginnote{\footnotesize\color{gray!50}\textsf{f.\,#1}}\ignorespaces}

% The modernised reading's coarser counterpart to \page{}: which views a
% section is drawn from.
\newcommand{\pagerange}[2]{%
  \marginnote{\footnotesize\color{gray!70}\textsf{v.\,#1--#2}}%
  \ignorespaces}

% Illegible: never guessed.
\newcommand{\ill}{\textcolor{red!60!black}{[\,\ldots\,]}}

% A doubtful reading: the word is given, and flagged as uncertain.
\newcommand{\uncertain}[1]{\uline{#1}}

% An editorial addition — a missing letter, an implied word.
\newcommand{\add}[1]{\textcolor{blue!50!black}{[#1]}}

% Struck out by the writer. What was crossed out is part of the document.
\newcommand{\struck}[1]{\sout{#1}}

% The transcriber's note, on the state of the page or the mathematics.
\newcommand{\note}[1]{\par\smallskip{\small\color{gray!60!black}\textit{[#1]}}\par\smallskip}

% A marginal note of hers, distinct from the body of the text.
\newcommand{\marginal}[1]{\par\smallskip\noindent\hspace{1em}%
  {\small\color{gray!60!black}#1}\par\smallskip}

% --- Title ------------------------------------------------------------------
% The title block is French, like the documents themselves.

\AddToHook{shipout/background}{%
  \ifx\grothWatermark\empty\else
    \begin{tikzpicture}[remember picture, overlay]
      \node[rotate=52, scale=3.4, align=center, text=gray!18,
            font=\sffamily\bfseries]
        at (current page.center) {\grothWatermark};
    \end{tikzpicture}%
  \fi}

\AtBeginDocument{%
  \begin{center}
    {\large\bfseries Manuscrits de mathématiciens (Gallica) ---
      \ifx\grothShelfmark\empty\grothFolder\else\grothShelfmark\fi}\\[2pt]
    {\itshape\grothFoldertitle}\\[4pt]
    {\small vues \grothFirst--\grothLast
      \ifx\grothBatch\empty\else\ (lot \grothBatch)\fi}\\[2pt]
    \ifx\grothDating\empty\else
      {\small\color{gray!60!black}Datation du catalogue : \grothDating}\\[2pt]
    \fi
    \ifx\grothWatermark\empty\else
      {\small\bfseries\color{red!45!gray}\grothWatermark}\\[2pt]
    \fi
    {\footnotesize\color{gray!60!black}Transcription établie d'après les images IIIF de Gallica.
     Source gallica.bnf.fr / Bibliothèque nationale de France.\\
     Les lectures incertaines sont soulignées, les illisibles notées [\,\ldots\,],
     les ajouts entre crochets ; \textsf{v.} numérote les vues de Gallica,
     \textsf{f.} la foliotation au crayon de la BnF.}
  \end{center}
  \vspace{1em}}

\endinput


\folder{fr-9115}
\shelfmark{Français 9115}
\batch{15}
\pages{281}{300}
\dating{1801-1900}
\watermark{Lecture automatique\\non vérifiée}
\foldertitle{Papiers de Sophie GERMAIN. — Recueil de dissertations et problèmes mathématiques et physiques.}

\begin{document}

\note{Numérotation des feuillets. Chaque feuillet écrit porte deux nombres,
tous deux semble-t-il à l'encre, et les versos blancs n'en portent aucun :
au milieu de la tête, souligné, de sa main, 30 (vue 283) puis 31 à 38 aux
vues 285 à 299 ; en haut à droite, d'une plume plus large et d'une écriture
plus ronde, 138 (vue 283) puis 139 à 146 aux vues 285 à 299, un par
feuillet et sans lacune. Cette seconde série occupe la place de la
foliotation de la Bibliothèque, mais elle est tracée à la plume et non au
crayon gris, et elle prolonge la série à l'encre relevée dans les lots
précédents de ce volume ; aucune foliotation au crayon n'a été lue ici, et
aucun « f. » n'est donc porté. Tous ces nombres ont été lus sur la vue ;
aucun n'est déduit. La vue 281 reprend le feuillet de la vue 280, béquet
relevé ; les vues 284, 286, 288, 290, 292, 294, 296, 298 et 300 sont les
versos blancs des feuillets qui les précèdent et ne sont pas transcrites.}

\page{281}

autre valeur de $x$ que \struck{$x$}, considérer la lame primitive
comme séparée en deux autres lames divisée au point
d'application du stilet \struck{et que parconséquent}
De sorte qu'il semble que ces formules indiquent
des conclusions qui ne sont pourtant \add{pas} admissibles \struck{que}
pour \struck{le cas ou}
\note{Cette vue photographie une seconde fois le feuillet de la vue 280,
le grand béquet rabattu vers le haut : on y lit les six lignes du bas de
la page, que le béquet cachait, et rien d'autre. Le haut du feuillet est
couvert par le béquet, dont le texte se lit à l'envers par transparence
et se rétablit en retournant la vue ; c'est le texte qui se lit à
l'endroit à la vue 280, et il n'est pas repris ici. Dans la marge de
gauche apparaissent les trois dernières lignes d'une longue note
marginale du même feuillet, elle aussi entière à la vue 280 ; elle n'est
pas reprise non plus. Tout le passage transcrit est barré d'une grande
croix à la plume ; un long trait horizontal passe en outre sur
« admissibles » et sur « le cas ou ». Le caractère biffé après « que »,
à la première ligne, est pris dans la croix et pourrait être un chiffre.
Aucun nombre n'est lisible en tête de cette vue.}

\page{283}

de sorte que pour une valeur quelconque de $x$ on peut \struck{la} toujours,
dans ce cas des extrémités appuiés, considérer la lame primitive
comme séparé en deux autres lames qui auroient egalement
leurs deux extrémités appuiés.

\marginal{mais ensuite}
\struck{Et} Pour les cas ou les extremités sont l'une et l'autre
libres ou fixés, on l'a vu \struck{que \ill{}} \add{parceque, comme} la lame primitive
ne peut être \struck{cassé} partagée en deux autres que \struck{pour} la
seule valeur $x=\frac{1}{2}$. \struck{mais} \add{la seule chose que montrent les équations du §.16 et que} pour une valeur quelconque de
$x$ \struck{ces équations montrent}, que si il y a un point d'appui
a la distance $ax$ d'une des extremités, il \struck{se} formera un
autre nœud a la même distance $ax$ de l'autre extremité, car ces
équations sont fonctions homogènes des quantités $ax$ et $a-ax$.

19 Je crois que l'on peut \struck{établir} \add{parvenir a} toutes ces \struck{résultats} \add{conséquences} sans avoir
recours a \struck{cette} \add{la} méthode \add{dont il s'agit ici} \struck{et} parconséquent sans employer la
considération des équations citées

En effet en se bornant a la seule considération de
l'état des extrémités \struck{\ill{}} \add{appuiés} \add{cas où} \struck{\ill{} de}
\add{que je suppose d'abord l'une \struck{et} l'autre appuiés}
\struck{l'ordonnée pour} on déduit des valeurs donnés par les cas et
que l'ordonnée pour une valeur quelconque de $x$ est
exprimée par l'équation $y = \dots \; \text{sin}\,x\omega$ de sorte que si par
une cause quelconque la courbe coupe l'axe dans le point
dont $y$ est l'ordonnée, ce qui arrive pour les nœuds,
on aura $\text{sin}\,x\omega = 0$. D'où on déduira egalement $\text{sin}\,2x\omega = \&c.$

Ensuite pour les extremités libres on a (Voyez ces premier)
lorsque la courbe coupe l'axe pour une valeur quelconque de $x$
\[ e^{x\omega} \pm e^{(1-x)\omega} + (1 \mp e^{\omega})\,\text{sin}\,x\omega + (1 \pm e^{\omega})\,\text{cos}\,x\omega = 0 \]
et en adoptant d'abord les signes inférieurs qui ont lieu lorsque
$\text{sin}\,\omega$ \struck{est positif}, qui est le cas ou il y a un nombre
impair on a d'abord en faisant $x=\frac{1}{2}$
\note{En tête, au milieu, « 30 » souligné, et en haut à droite « 138 », tous
deux à l'encre ; le second d'une plume plus large. Le paragraphe
« Je crois… » est précédé d'un « 19 » empâté, écrit dans la ligne : les
paragraphes de ces feuillets sont numérotés. Deux groupes biffés n'ont
pas été lus : après « on l'a vu », trois mots dont le dernier se termine
par « -ait » ; et, à la ligne « l'état des extrémités », deux groupes,
le second suivi de « de ». Le « $y = \dots$ » laisse en blanc le facteur
de tête, comme ailleurs dans ces papiers. Les longs traits horizontaux
qui passent sur « extremités » et « sont » sont ses traits de liaison et
ne sont pas donnés pour des ratures. Le verso de ce feuillet (vue 284)
est blanc.}

\page{285}

\[ (1+e^{\omega})\,\text{sin}\,\tfrac{1}{2}\omega + (1-e^{\omega})\,\text{cos}\,\tfrac{1}{2}\omega = 0 \quad\text{ou}\quad \text{tang}\,\tfrac{1}{2}\omega = \frac{e^{\omega}-1}{e^{\omega}+1} \]
équation qui résulte des conditions $\text{cos}\,\omega = \frac{2e^{\omega}}{e^{2\omega}+1}$ et $\text{sin}\,\omega = \frac{e^{2\omega}-1}{e^{2\omega}+1}$
comme on l'a déjà vu

reprenant ensuite l'équation générale
\[ e^{x\omega} - e^{(1-x)\omega} + (1+e^{\omega})\,\text{sin}\,x\omega + (1-e^{\omega})\,\text{cos}\,x\omega = 0 \]
\struck{on voit que l'on peut y ajouter}, sans qu'elle \struck{\uncertain{cesse d'être}}
et la retranchant de $(1+e^{\omega})\,\text{sin}\,\frac{1}{2}\omega + (1-e^{\omega})\,\text{cos}\,\frac{1}{2}\omega = 0$ multipliée par $2\,\text{cos}\,\frac{(1-2x)\omega}{2}$
c'est a dire de $2\{(1+e^{\omega})\,\text{sin}\,\frac{1}{2}\omega + (1-e^{\omega})\,\text{cos}\,\frac{1}{2}\omega\}\,\text{cos}\,\frac{(1-2x)\omega}{2}$
\[ = (1+e^{\omega})\{\text{sin}\,x\omega + \text{sin}(1-x)\omega\} + (1-e^{\omega})\{\text{cos}\,x\omega + \text{cos}(1-x)\omega\} = 0 \]
on trouve
\[ e^{(1-x)\omega} - e^{x\omega} + (1+e^{\omega})\{\text{sin}(1-x)\omega + \text{sin}\,x\omega - \text{sin}\,x\omega\} \]
\[ + (1-e^{\omega})\{\text{cos}(1-x)\omega + \text{cos}\,x\omega - \text{cos}\,x\omega\} \]
\[ = e^{(1-x)\omega} - e^{x\omega} + (1+e^{\omega})\,\text{sin}(1-x)\omega + (1-e^{\omega})\,\text{cos}(1-x)\omega = 0 \]
Ainsi l'ordonnée est encore zéro lorsque l'on change $x$ en
$x-1$.

Demême aussi pour le cas ou il y a un nombre pair
de nœuds, si on \struck{ajoutée} retranche \add{de} l'équation
\[ e^{x\omega} + e^{(1-x)\omega} + (1-e^{\omega})\,\text{sin}\,x\omega + (1+e^{\omega})\,\text{cos}\,x\omega = 0 \]
cette autre équation $(e^{\omega}+1)\,\text{sin}\,\frac{1}{2}\omega - (1-e^{\omega})\,\text{cos}\,\frac{1}{2}\omega = 0$ tirée des
conditions $\text{cos}\,\omega = \frac{2e^{\omega}}{e^{2\omega}+1}$, $\text{sin}\,\omega = -\frac{e^{2\omega}-1}{e^{2\omega}+1}$, après l'avoir multipliée
par $2\,\text{sin}\,\frac{(1-2x)\omega}{2}$; ce qui donne
\[ \{(e^{\omega}+1).2\,\text{sin}\,\tfrac{1}{2}\omega - (1-e^{\omega}).2\,\text{cos}\,\tfrac{1}{2}\omega\}\,\text{sin}\,\frac{(1-2x)\omega}{2} \]
\[ = (e^{\omega}+1)(\text{cos}\,x\omega - \text{cos}(1-x)\omega) \ill{} (1-e^{\omega})(\text{sin}(1-x)\omega - \text{sin}\,x\omega) \]
on trouve
\[ e^{x\omega} + e^{(1-x)\omega} + (1-e^{\omega})\{\text{sin}(1-x)\omega - \text{sin}\,x\omega + \text{sin}\,x\omega\} \]
\[ - (1+e^{\omega})\{\text{cos}\,x\omega - \text{cos}(1-x)\omega - \text{cos}\,x\omega\} \]
\[ = e^{(1-x)\omega} + e^{x\omega} + (1-e^{\omega})\,\text{sin}(1-x)\omega + (1+e^{\omega})\,\text{cos}(1-x)\omega = 0 \]
c'est a dire que l'ordonnée est egalement zéro lorsque l'on change
$x$ en $1-x$.

Les résultats sont parfaitement semblables pour les cas ou
les deux extrémités sont fixés, car la seule différence
\note{En tête, au milieu, « 31 » souligné ; en haut à droite « 139 », à
l'encre. « lorsque l'on change $x$ en $x-1$ » est écrit ainsi sur le
feuillet ; la même conclusion est écrite « en $1-x$ » quinze lignes plus
bas, et le calcul donne $1-x$ : la leçon est laissée telle quelle. Le
signe qui joint les deux produits de l'avant-dernier groupe est sous un
gros pâté et n'est pas lu. Un pâté couvre aussi le « 1 » de
« $\text{cos}(1-x)\omega$ » à la troisième équation avant la fin, et le
« 1 » de « $(1-e^{\omega})$ » dans « cette autre équation ». « si on
retranche de l'équation » porte trois états superposés : « ajoutée »
biffé au-dessus de la ligne, « retranche » biffé sur la ligne et repris
« tranche de » au-dessous, un « a » au-dessus de « de ». Le verso
(vue 286) est blanc.}

\page{287}

entre ce cas et le précédant est qu'on a ici
\[ e^{x\omega} \pm e^{(1-x)\omega} - \{(1 \mp e^{\omega})\,\text{sin}\,x\omega + (1 \pm e^{\omega})\,\text{cos}\,x\omega\} = 0 \]
au lieu de
\[ e^{x\omega} \pm e^{(1-x)\omega} + (1 \mp e^{\omega})\,\text{sin}\,x\omega + (1 \pm e^{\omega})\,\text{cos}\,x\omega = 0 \]

\marginal{que les équations finales semblent \struck{présenter} \add{indiquer}
des conclusions qui ne sont pourtant pas admissibles, puisqu'elles sont
inconcilliables avec les hypothèses subsistantes sur l'état des extrèmités
\struck{et} \add{après} avoir prouvé que le petit nombre de conséquences
applicables que l'on peut tirer de ces équations \struck{peuvent} etres
atteintes \struck{d'une manière} \struck{independantes des mêmes équations}
par des considérations fort simples et independantes des mêmes equations}

20 Après avoir remarqué $*$ que le petit nombre de conclusions
qui peuvent, ayant egard a l'état des extremités, être
tirée des équations trouvées par la methode d'Euler
peuvent etre déduits avec beaucoup plus de simplicité
de l'expression de l'ordonnée convenable a chaque
cas et supposée égale a zéro dans le point
ou il y a un nœud; soit que ce nœud soit l'effet
immédiat de l'application d'un stilet soit
qu'il se forme par une \struck{autre} cause, il étoit
naturel d'examiner les \struck{hypothèses} \add{la methode qui a servi a les trouver} sur lesquelles
\struck{cette méthode est appuiée} \add{au moyen de laquelle elles ont été trouvées}

On peut voir N\textsuperscript{o} 47 \add{dans le \uncertain{quinzième} du memoire} \add{que l'auteur} \struck{que d'Euler}, après avoir admis que
le stilet soit simplement appuié de manière a ce que le
le mouvement puisse se communiquer de l'une a l'autre
des deux parties qu'il sépare suppose cependant que
l'action de ce stilet interrompe la continuité de la
courbe de sorte que les deux portions \add{suivant cette supposition} de \struck{cette} \add{cette} courbe
seroient exprimés par des équations qui \struck{différent} \add{seroient} seulement
par les valeurs des coefficiens \add{$\beta$, $\gamma$ etc,} et que c'est d'après cette
hypothese qu'il a établi les équations I II III
IV

\struck{Indépendamment} \add{sans avoir recours au} de tout calcul, il peut d'abord paroitre
étonnant \struck{qu'Euler admette qu'elle} \add{qu'Euler suppose que les deux portions séparés par l'application du stilet n'appartiennent plus a une seule et même courbe $+$} l'application du
stilet produisant un simple point d'appui, puisse
\marginal{que le mouvement dont il s'agit etant du nombre des mouvement
reguliers de la lame puisque le seul effet de l'application du stilet est
de déterminer la formation d'un certain nombre de nœuds c'est a dire}
de donner lieu a un des mouvement \struck{réguliers déterminés} \add{réguliers dont il s'agit \uncertain{ici en question}} dans
\note{En tête, au milieu, « 32 » souligné ; en haut à droite « 140 », à
l'encre. Le paragraphe est numéroté « 20 » dans la ligne. Le feuillet
porte deux notes de sa main dans la marge de gauche, données ici avant
et après la ligne qu'elles accompagnent : la première répond à
l'astérisque de « Après avoir remarqué * », la seconde à la croix « + »
de l'addition sur Euler ; la seconde est barrée d'un long trait
oblique. Les chiffres romains « I II III IV » qui closent le paragraphe
sont les numéros des équations d'Euler, écrits d'une plume plus large.
« N\textsuperscript{o} 47 » renvoie au paragraphe 47 du mémoire d'Euler,
celui qui est copié aux vues 192 et suivantes. Ce feuillet est repris
de part en part : les mots donnés en addition sont écrits au-dessus ou
au-dessous de la ligne, sans signe d'insertion. Plusieurs longs traits
horizontaux passent sur des mots qui font sens et ne sont pas donnés
pour des ratures. Le verso (vue 288) est blanc.}

\page{289}

\marginal{l'examen des six cas relatifs a \add{a l'exclusion des autres}
l'état des extremités, et les mouvement dont il s'agit se fesant sans
interruption dans les courbes auxquelles ils donnent lieu, Euler suppose
pourtant que les deux portions séparés par le stilet n'appartienne plus
a une seule et même courbe.}

donner lieu a la discontinuité dela courbe. en effet
on a vu dans l'examen des six cas relatifs a l'état des
extremités, que l'on parvient a des équations qui déterminent
les valeurs convenables de l'angle $\omega$ et que ces differentes
valeurs sont relatifs aux differens mouvemens reguliers dont
les lames auxquelles elles s'appliquent sont susceptibles
mais le mouvement modifié par l'application du
stilet simplement appuié est \add{encore} d'après l'hypothèse,
N\textsuperscript{o} 47, un mouvement regulier des mêmes lames, il
doit donc etre compris au nombre de ceux
qui sont déterminés par les conditions données
dans l'examen des six cas cités et qui n'admettent
pas de discontinuité dans la courbe a laquelle ils
donnent lieu. On est \struck{ainsi menée} a conclure que les coefficiens
$\alpha'$, $\beta'$, $\gamma'$ et $\delta'$ sont égaux chacun a son corrispondant parmi les
coefficiens $\alpha$, $\beta$, $\gamma$, et $\delta$. $+$

\marginal{$+$ car le seul effet de cette application est de déterminer la
formation d'un certain nombre de nœuds et parconséquent de donner lieu
a de certains mouvemens reguliers compris au nombre de ceux qui ont été
traités dans l'examen des six cas relatifs a l'état des extremités
\struck{a l'exclus} \struck{et} les mouvemens dont il s'agit s'éxécutent
suivant une courbe continue. \struck{et} ces réflexions mènent donc
\struck{cependant je}}

\add{Cependant} \struck{Ce résultat} Je n'ose admettre ce resultat contraire a
l'opinion d'Euler. sur une hypothèse \struck{qui} \add{par laquelle il} semble avoir
ou \struck{devoir} ouvrir la carrière a l'examen d'une foule de
cas dont il donne \struck{l'examen d'un seul comme un}
\struck{exemple qui doit guider} un exemple comme pour guider
dans des recherches semblables, sans le soumettre
a un examen plus approfondi.
c'est pourquoi je me propose de démontrer
la proposition suivante :

\marginal{Cependant je n'ose admettre une opinion si contraire a celle
d'Euler sans la soumettre a un examen plus approfondi c'est pourquoi je
me propose de demontrer la proposition suivante}
\note{En tête, au milieu, « 33 » souligné ; en haut à droite « 141 », à
l'encre. Le corps de la page est barré d'une grande croix à la plume, et
la première note marginale d'un long trait oblique ; les mots biffés au
trait bref sont seuls donnés ici comme ratures. La troisième note
marginale reprend, presque mot pour mot, les trois dernières lignes du
corps : elle en est la mise au net. Les lettres grecques sont celles
d'Euler : le $\gamma$ est tracé comme un « v », ainsi que l'a déjà noté
une passe antérieure sur la copie latine de ce volume.
« N\textsuperscript{o} 47 » renvoie encore au paragraphe 47 du mémoire
d'Euler. Le bas du feuillet est blanc ; le verso (vue 290) l'est aussi.}

\page{291}

21 Toutes les fois que la condition de l'application d'un stilet
simplement appuié dans un point quelconque de l'étendue
d'une lame elastique vibrante: laisse subsister une
hypothese \struck{quelconque} \add{déterminée} sur l'état des extrimités de la
meme lame. les coëfficiens $\alpha$, $\beta$, $\gamma$, $\delta$ et $\alpha'$, $\beta'$, $\gamma'$, $\delta'$ qui
entrent dans l'éxpression des ordonnées des deux parties
de la courbe séparées par ce point d'appui sont égaux
chacun a son corrspondant

Je vais commencer par prouver qu'\add{en général} il suffit que l'un de ces coëfficiens
soit égal a son correspondant pour que la même égalité ait
lieu entre tous les autres. De sorte que si on a $\alpha = \alpha'$ il
en resultera nécessairement $\beta = \beta'$, $\gamma' = \gamma$, \&c.

Je suppose donc $\alpha = \alpha'$, il resulte alors de l'équation \struck{\ill{}}
$\alpha e^{x\omega} + \beta e^{-x\omega} = \alpha' e^{x\omega} + \beta' e^{-x\omega}$ \struck{\ill{}} que $\beta = \beta'$ dans cette hypothese
les équations III et IV se réduisent a
\[ \text{III.}\ (\gamma-\gamma')\,\text{cos}\,x\omega - (\delta-\delta')\,\text{sin}\,x\omega = 0, \quad \text{IV}\ (\gamma-\gamma')\,\text{sin}\,x\omega + (\delta-\delta')\,\text{cos}\,x\omega = 0 \]
la première multipliée par $\text{cos}\,x\omega$ et ajoutée a la seconde multipliée
par $\text{sin}\,x\omega$ donne $(\gamma-\gamma')(\text{sin}^2 x\omega + \text{cos}^2 x\omega) = 0$ c'est adire $\gamma = \gamma'$
Ainsi \struck{et} on a III $(\delta-\delta')\,\text{sin}\,x\omega = 0$ et IV $(\delta-\delta')\,\text{cos}\,x\omega = 0$ et
parséquent $\delta = \delta'$

Il ne s'agit donc plus que de prouver qu'en effet la
condition $\alpha = \alpha'$ est necessaire, je choisis d'abord le cas
ou les deux extremités sont supposées libres.

On a alors \struck{l'équation}
\[ \frac{\alpha[(e^{x\omega}+e^{-x\omega})\,\text{sin}\,x\omega - (e^{x\omega}-e^{-x\omega})\,\text{cos}\,x\omega]}{\text{sin}\,x\omega - \text{cos}\,x\omega - e^{x\omega}} \]
\[ = \frac{\alpha'[(e^{(1-x)\omega}+e^{(x-1)\omega})\,\text{sin}(1-x)\omega - (e^{(1-x)\omega}-e^{(x-1)\omega})\,\text{cos}(1-x)\omega]}{e^{-x\omega}+e^{-\omega}(\text{cos}(1-x)\omega + \text{sin}(1-x)\omega)} \]
ainsi il suffit de \struck{pro} montrer que l'équation
\[ \{e^{-x\omega}+e^{-\omega}(\text{cos}(1-x)\omega+\text{sin}(1-x)\omega)\}\{(e^{x\omega}+e^{-x\omega})\,\text{sin}\,x\omega - (e^{x\omega}-e^{-x\omega})\,\text{cos}\,x\omega\} \]
\[ = \{\text{sin}\,x\omega - \text{cos}\,x\omega - e^{-x\omega}\}\{(e^{(1-x)\omega}+e^{(x-1)})\,\text{sin}(1-x)\omega - (e^{(1-x)\omega}-e^{(x-1)\omega})\,\text{cos}(1-x)\omega\} \]
que \struck{l'équa} a lieu d'ellemême sous les conditions indiqués.
\note{En tête, au milieu, « 34 » souligné ; en haut à droite « 142 », à
l'encre. Le paragraphe porte le numéro 21, écrit dans la ligne et
empâté. « parséquent » est écrit ainsi, contraction de « par
conséquent », qu'elle écrit ailleurs « parconséquent ». Le $\gamma$ est
tracé comme un « v ». Le dénominateur de la première fraction porte
« $e^{x\omega}$ » sans signe à l'exposant, alors que le même facteur est
écrit « $e^{-x\omega}$ » quatre lignes plus bas ; les deux leçons sont
laissées telles. Dans le second membre de la dernière équation, un
« $e^{(x-1)}$ » est écrit sans $\omega$ à l'exposant, à la différence
des trois autres. Deux groupes courts biffés, l'un après « de
l'équation », l'autre après le second membre de l'équation en $\alpha$
et $\beta$, ne sont pas lus. Le verso (vue 292) est blanc.}

\page{293}

On a par le developpement:
\[ (e^{(x-1)\omega} + e^{-(x+1)\omega})\{\text{cos}(1-x)\omega\,\text{sin}\,x\omega + \text{sin}(1-x)\omega\,\text{sin}\,x\omega\} - (e^{(x-1)\omega} - e^{-(x+1)\omega})\{\text{cos}(1-x)\omega\,\text{cos}\,x\omega + \text{sin}(1-x)\omega\,\text{cos}\,x\omega\} \]
\[ + e^{-x\omega}\{(e^{(1-x)\omega} + e^{(x-1)\omega})\,\text{sin}(1-x)\omega - (e^{(1-x)\omega} - e^{(x-1)\omega})\,\text{cos}(1-x)\omega\} \]
\[ = (e^{(1-x)\omega} + e^{(x-1)\omega})\{\text{sin}(1-x)\omega\,\text{sin}\,x\omega - \text{sin}(1-x)\omega\,\text{cos}\,x\omega\} - (e^{(1-x)\omega} - e^{(x-1)\omega})\{\text{cos}(1-x)\omega\,\text{sin}\,x\omega - \text{cos}(1-x)\omega\,\text{cos}\,x\omega\} \]
\[ - e^{-x\omega}\{(e^{x\omega} + e^{-x\omega})\,\text{sin}\,x\omega - (e^{x\omega} - e^{-x\omega})\,\text{cos}\,x\omega\} \]
ou
\[ e^{(x-1)\omega}\{\text{cos}(1-x)\omega\,\text{sin}\,x\omega + \text{sin}(1-x)\omega\,\text{sin}\,x\omega - \text{cos}(1-x)\omega\,\text{cos}\,x\omega - \text{sin}(1-x)\omega\,\text{cos}\,x\omega\} \]
\[ + e^{-(1+x)\omega}\{\text{cos}(1-x)\omega\,\text{sin}\,x\omega + \text{sin}(1-x)\omega\,\text{sin}\,x\omega + \text{cos}(1-x)\omega\,\text{cos}\,x\omega + \text{sin}(1-x)\omega\,\text{cos}\,x\omega\} \]
\[ + e^{-x\omega}\{(e^{(1-x)\omega} + e^{(x-1)\omega})\,\text{sin}(1-x)\omega - (e^{(1-x)\omega} - e^{(x-1)\omega})\,\text{cos}(1-x)\omega + (e^{x\omega} + e^{-x\omega})\,\text{sin}\,x\omega - (e^{x\omega} - e^{-x\omega})\,\text{cos}\,x\omega\} \]
\[ = e^{(x-1)\omega}\{\text{sin}(1-x)\omega\,\text{sin}\,x\omega - \text{sin}(1-x)\omega\,\text{cos}\,x\omega + \text{cos}(1-x)\omega\,\text{sin}\,x\omega - \text{cos}(1-x)\omega\,\text{cos}\,x\omega\} \]
\[ + e^{(1-x)\omega}\{\text{sin}(1-x)\omega\,\text{sin}\,x\omega - \text{sin}(1-x)\omega\,\text{cos}\,x\omega - \text{cos}(1-x)\omega\,\text{sin}\,x\omega + \text{cos}(1-x)\omega\,\text{cos}\,x\omega\} \]
en effaceant ce qui se détruit et multipliant par $e^{-x\omega}$ cette
équation devient :
\[ e^{-\omega}\{\text{cos}(1-x)\omega\,\text{sin}\,x\omega + \text{sin}(1-x)\omega\,\text{sin}\,x\omega + \text{cos}(1-x)\omega\,\text{cos}\,x\omega + \text{sin}(1-x)\omega\,\text{cos}\,x\omega\} \]
\[ + (e^{(1-x)\omega} + e^{(x-1)\omega})\,\text{sin}(1-x)\omega - (e^{(1-x)\omega} - e^{(x-1)\omega})\,\text{cos}(1-x)\omega + (e^{x\omega} + e^{-x\omega})\,\text{sin}\,x\omega - (e^{x\omega} - e^{-x\omega})\,\text{cos}\,x\omega \]
\[ = e^{\omega}\{\text{sin}(1-x)\omega\,\text{sin}\,x\omega - \text{sin}(1-x)\omega\,\text{cos}\,x\omega - \text{cos}(1-x)\omega\,\text{sin}\,x\omega + \text{cos}(1-x)\omega\,\text{cos}\,x\omega\} \]
et elle peut etre mise sous la forme
\[ e^{-\omega}\{\text{cos}(1-x)\omega \ill{} \text{sin}(1-x)\omega\}\{\text{cos}\,x\omega + \text{sin}\,x\omega\} \]
\[ + e^{(1-x)\omega}\{\text{sin}(1-x)\omega - \text{cos}(1-x)\omega\} + e^{(x-1)\omega}\{\text{sin}(1-x)\omega + \text{cos}(1-x)\omega\} + e^{x\omega}\{\text{sin}\,x\omega - \text{cos}\,x\omega\} + e^{-x\omega}\{\text{sin}\,x\omega + \text{cos}\,x\omega\} \]
\[ = e^{\omega}\{\text{sin}(1-x)\omega - \text{cos}(1-x)\omega\}\{\text{sin}\,x\omega - \text{cos}\,x\omega\} \]
je suppose d'abord $\text{sin}\,\omega$ positif, on a pour ce cas
\[ e^{x\omega} - e^{(1-x)\omega} + (1+e^{\omega})\,\text{sin}\,x\omega + (1-e^{\omega})\,\text{cos}\,x\omega = 0 \quad \text{d'ou il resulte} \]
\[ e^{(1-x)\omega} - e^{x\omega} + (1+e^{\omega})\,\text{sin}(1-x)\omega + (1-e^{\omega})\,\text{cos}(1-x)\omega = 0 \]
et parconséquent
\[ \{\text{sin}\,x\omega + \text{cos}\,x\omega\}\{\text{sin}(1-x)\omega + \text{cos}(1-x)\omega\} = \{e^{x\omega} - e^{(1-x)\omega} + e^{\omega}(\text{sin}\,x\omega - \text{cos}\,x\omega)\}\{e^{(1-x)\omega} - e^{x\omega} + e^{\omega}(\text{sin}(1-x)\omega - \text{cos}(1-x)\omega)\} \]
\[ = e^{\omega} - e^{2(1-x)\omega} - e^{2x\omega} + e^{\omega} + e^{\omega}\{e^{x\omega}(\text{sin}(1-x)\omega - \text{cos}(1-x)\omega) - e^{(1-x)\omega}(\text{sin}(1-x)\omega - \text{cos}(1-x)\omega)\} \]
\[ + e^{\omega}\{e^{(1-x)\omega}(\text{sin}\,x\omega - \text{cos}\,x\omega) - e^{x\omega}(\text{sin}\,x\omega - \text{cos}\,x\omega)\} \]
\[ + e^{2\omega}\{\text{sin}\,x\omega - \text{cos}\,x\omega\}\{\text{sin}(1-x)\omega - \text{cos}(1-x)\omega\} \]
en mettant cette valeur dans l'equation cidessus elle devient:
\note{En tête, au milieu, « 35 » souligné ; en haut à droite « 143 », à
l'encre. La page est tout entière de calcul : aucun paragraphe numéroté.
Dans « elle peut etre mise sous la forme », le signe qui joint
« $\text{cos}(1-x)\omega$ » et « $\text{sin}(1-x)\omega$ » est sous un
gros pâté et n'est pas lu. La ligne qui suit ce premier facteur est
d'abord écrite comme la ligne correspondante plus haut, puis biffée de
traits en croix et récrite au-dessous sous la forme des quatre
exponentielles : c'est cette seconde rédaction qui est donnée ici. Le
feuillet s'arrête sur « elle devient: » ; la suite est à la vue 295. Le
verso (vue 294) est blanc.}

\page{295}

\[ 2 - e^{(1-2x)\omega} - e^{(2x-1)\omega} + e^{x\omega}\{\text{sin}(1-x)\omega - \text{cos}(1-x)\omega\} - e^{(1-x)\omega}\{\text{sin}(1-x)\omega - \text{cos}(1-x)\omega\} \]
\[ + e^{(1-x)\omega}\{\text{sin}\,x\omega - \text{cos}\,x\omega\} - e^{x\omega}\{\text{sin}\,x\omega - \text{cos}\,x\omega\} \]
\[ + e^{\omega}\{\text{sin}\,x\omega - \text{cos}\,x\omega\}\{\text{sin}(1-x)\omega - \text{cos}(1-x)\omega\} \]
\[ + e^{(1-x)\omega}\{\text{sin}(1-x)\omega - \text{cos}(1-x)\omega\} + e^{(x-1)\omega}\{\text{sin}(1-x)\omega + \text{cos}(1-x)\omega\} + e^{x\omega}\{\text{sin}\,x\omega - \text{cos}\,x\omega\} + e^{-x\omega}\{\text{sin}\,x\omega + \text{cos}\,x\omega\} \]
\[ = e^{\omega}\{\text{sin}\,x\omega - \text{cos}\,x\omega\}\{\text{sin}(1-x)\omega - \text{cos}(1-x)\omega\} \]
et en effaceant ce qui se détruit
\[ 2 - e^{(1-2x)\omega} - e^{(2x-1)\omega} + e^{x\omega}\{\text{sin}(1-x)\omega - \text{cos}(1-x)\omega\} + e^{(x-1)\omega}\{\text{sin}(1-x)\omega + \text{cos}(1-x)\omega\} \]
\[ + e^{(1-x)\omega}\{\text{sin}\,x\omega - \text{cos}\,x\omega\} + e^{-x\omega}\{\text{sin}\,x\omega + \text{cos}\,x\omega\} = 0 \]
en reprenant l'équation
\[ e^{x\omega} - e^{(1-x)\omega} + (1+e^{\omega})\,\text{sin}\,x\omega + (1-e^{\omega})\,\text{cos}\,x\omega = 0 \]
et la divisant par $e^{x\omega}$ elle donne
\[ 1 - e^{(1-2x)\omega} + e^{-x\omega}\{\text{sin}\,x\omega + \text{cos}\,x\omega\} + e^{(1-x)\omega}\{\text{sin}\,x\omega - \text{cos}\,x\omega\} = 0 \]
\uncertain{demême} on tiré de l'équation
\[ e^{(1-x)\omega} - e^{x\omega} + (1+e^{\omega})\,\text{sin}(1-x)\omega + (1-e^{\omega})\,\text{cos}(1-x)\omega = 0 \]
en la divisant par $e^{(1-x)\omega}$
\[ 1 - e^{(2x-1)\omega} + e^{(x-1)\omega}\{\text{sin}(1-x)\omega + \text{cos}(1-x)\omega\} + e^{x\omega}\{\text{sin}(1-x)\omega - \text{cos}(1-x)\omega\} = 0 \]
Enfin de l'addition de ces deux équations il resulte
\[ 2 - e^{(2x-1)\omega} - e^{(2x-1)\omega} + e^{-x\omega}\{\text{sin}\,x\omega + \text{cos}\,x\omega\} + e^{(1-x)\omega}\{\text{sin}\,x\omega - \text{cos}\,x\omega\} \]
\[ + e^{(x-1)\omega}\{\text{sin}(1-x)\omega + \text{cos}(1-x)\omega\} + e^{x\omega}\{\text{sin}(1-x)\omega - \text{cos}(1-x)\omega\} = 0 \]
ainsi l'équation est demontrée pour le cas ou $\text{sin}\,\omega$ est positif

je suppose apresent $\text{sin}\,\omega$ negatif on aura
\[ e^{x\omega} + e^{(1-x)\omega} + (1-e^{\omega})\,\text{sin}\,x\omega + (1+e^{\omega})\,\text{cos}\,x\omega = 0 \]
\[ e^{x\omega} + e^{(1-x)\omega} + (1-e^{\omega})\,\text{sin}(1-x)\omega + (1+e^{\omega})\,\text{cos}(1-x)\omega = 0 \]
\[ (\text{sin}\,x\omega + \text{cos}\,x\omega)(\text{sin}(1-x)\omega + \text{cos}(1-x)\omega) \]
\[ = e^{2\omega}(\text{sin}\,x\omega - \text{cos}\,x\omega)(\text{sin}(1-x)\omega - \text{cos}(1-x)\omega) - \{e^{x\omega} + e^{(1-x)\omega}\}\{\text{sin}\,x\omega - \text{cos}\,x\omega + \text{sin}(1-x)\omega - \text{cos}(1-x)\omega\}e^{\omega} \]
\[ + 2e^{\omega} + \struck{\ill{}} + e^{2x\omega} + e^{2(1-x)\omega} \]
mettant donc cette valeur dans l'équation
\[ e^{-\omega}\{\text{cos}(1-x)\omega + \text{sin}(1-x)\omega\}\{\text{cos}\,x\omega + \text{sin}\,x\omega\} \]
\[ \ill{} e^{(1-x)\omega}\{\text{sin}(1-x)\omega - \text{cos}(1-x)\omega\} + e^{(x-1)\omega}\{\text{sin}(1-x)\omega + \text{cos}(1-x)\omega\} + e^{x\omega}\{\text{sin}\,x\omega - \text{cos}\,x\omega\} + e^{-x\omega}\{\text{sin}\,x\omega - \text{cos}\,x\omega\} \]
\[ = e^{\omega}\{\text{sin}\,x\omega - \text{cos}\,x\omega\}\{\text{sin}(1-x)\omega - \text{cos}(1-x)\omega\} \]
elle devient :
\note{En tête, au milieu, « 36 » souligné ; en haut à droite « 144 », à
l'encre, dans un angle noirci par l'ombre du feuillet. Le premier terme
de l'équation qui suit « Enfin de l'addition de ces deux équations » est
écrit « $-e^{(2x-1)\omega} - e^{(2x-1)\omega}$ », le même exposant deux
fois, là où les deux équations additionnées donnent
« $-e^{(1-2x)\omega} - e^{(2x-1)\omega}$ » ; la leçon du feuillet est
gardée. De même, à l'avant-dernier groupe de la page, le dernier terme
est écrit « $e^{-x\omega}\{\text{sin}\,x\omega - \text{cos}\,x\omega\}$ »
alors que le même terme porte « $+\,\text{cos}\,x\omega$ » plus haut. Un
groupe biffé et repassé après « $+2e^{\omega}$ » n'est pas lu ; un gros
pâté couvre le signe qui ouvre la deuxième ligne du dernier groupe.
« demême » est lu avec doute : le mot est traversé d'un trait et son
milieu est empâté. Le verso (vue 296) est blanc.}

\page{297}

\[ e^{\omega}\{\text{sin}\,x\omega - \text{cos}\,x\omega\}\{\text{sin}(1-x)\omega - \text{cos}(1-x)\omega\} + 2 + e^{(2x-1)\omega} + e^{(1-2x)\omega} \]
\[ - e^{x\omega}\{\text{sin}\,x\omega - \text{cos}\,x\omega\} - e^{x\omega}\{\text{sin}(1-x)\omega - \text{cos}(1-x)\omega\} - e^{(1-x)\omega}\{\text{sin}\,x\omega - \text{cos}\,x\omega\} - e^{(1-x)\omega}\{\text{sin}(1-x)\omega - \text{cos}(1-x)\omega\} \]
\[ + e^{x\omega}\{\text{sin}\,x\omega - \text{cos}\,x\omega\} + e^{-x\omega}\{\text{sin}\,x\omega + \text{cos}\,x\omega\} + e^{(1-x)\omega}\{\text{sin}(1-x)\omega - \text{cos}(1-x)\omega\} + e^{(x-1)\omega}\{\text{sin}(1-x)\omega + \text{cos}(1-x)\omega\} \]
\[ = e^{\omega}\{\text{sin}\,x\omega - \text{cos}\,x\omega\}\{\text{sin}(1-x)\omega - \text{cos}(1-x)\omega\} \]
et en effaceant ce qui se détruit
\[ 2 + e^{(2x-1)\omega} + e^{(1-2x)\omega} + e^{-x\omega}\{(1-e^{\omega})\,\text{sin}\,x\omega + (1+e^{\omega})\,\text{cos}\,x\omega\} \]
\[ + e^{(x-1)\omega}\{(1-e^{\omega})\,\text{sin}(1-x)\omega + (1+e^{\omega})\,\text{cos}(1-x)\omega\} = 0 \]
On voit que comme dans le cas précédant cette équation
est la somme des deux équations
\[ e^{x\omega} + e^{(1-x)\omega} + (1-e^{\omega})\,\text{sin}\,x\omega + (1+e^{\omega})\,\text{cos}\,x\omega = 0 \]
\[ e^{x\omega} + e^{(1-x)\omega} + (1-e^{\omega})\,\text{sin}(1-x)\omega + (1+e^{\omega})\,\text{cos}(1-x)\omega = 0 \]
dont la première est multipliée par $e^{-x\omega}$ et la seconde
par $e^{(1-x)\omega}$.

22 Lorsque les équations sont l'une et l'autre \struck{appuiées}
fixées il est également vrai que les deux séries de coëfficiens
$\alpha$, $\beta$, $\gamma$, $\delta$ et $\alpha'$, $\beta'$, $\gamma'$, $\delta'$ sont égales entr'elles

En effet on a pour ce dernier cas
\[ \frac{\alpha}{\alpha'} = \frac{(e^{(1-x)\omega} + e^{(x-1)\omega})\,\text{sin}(1-x)\omega - (e^{(1-x)\omega} - e^{(x-1)\omega})\,\text{cos}(1-x)\omega}{(e^{x\omega} + e^{-x\omega})\,\text{sin}\,x\omega - (e^{x\omega} - e^{-x\omega})\,\text{cos}\,x\omega} \cdot \frac{\text{sin}\,x\omega - \text{cos}\,x\omega + e^{-x\omega}}{e^{-\omega}\{\text{sin}(1-x)\omega + \text{cos}(1-x)\omega\} - e^{-x\omega}} \]
au lieu de \struck{\ill{}}
\[ \frac{\alpha}{\alpha'} = \frac{(e^{(1-x)\omega} + e^{(x-1)\omega})\,\text{sin}(1-x)\omega - (e^{(1-x)\omega} - e^{(x-1)\omega})\,\text{cos}(1-x)\omega}{(e^{x\omega} + e^{-x\omega})\,\text{sin}\,x\omega - (e^{x\omega} - e^{-x\omega})\,\text{cos}\,x\omega} \cdot \frac{\text{sin}\,x\omega - \text{cos}\,x\omega - e^{-x\omega}}{e^{-\omega}\{\text{sin}(1-x)\omega + \text{cos}(1-x)\omega\} + e^{-x\omega}} \]
que l'on a trouvé pour le cas précédant, il est aisé \add{devoir} que la seule
Différence qui existe entre ces deux équations consiste dans le changement
de signe de la quantité $e^{-x\omega}$ et comme on a pour les
extremités \struck{appuiées} fixées
\[ e^{x\omega} \mp e^{(1-x)\omega} - \{(1 \pm e^{\omega})\,\text{sin}\,x\omega + (1 \mp e^{\omega})\,\text{cos}\,x\omega\} = 0 \]
au lieu de
\[ e^{x\omega} \mp e^{(1-x)\omega} + (1 \pm e^{\omega})\,\text{sin}\,x\omega + (1 \mp e^{\omega})\,\text{cos}\,x\omega = 0 \]
pour les extremités libres, il y a compensation entre les changemens
de signes et on parvient ici comme dans le cas précedant a
conclure $\alpha = \alpha'$.
\note{En tête, au milieu, « 37 » souligné ; en haut à droite « 145 », à
l'encre, dans un angle noirci. Le paragraphe porte le numéro 22, écrit
dans la ligne. À son premier mot elle écrit « Lorsque les équations sont
l'une et l'autre fixées » là où le sens demande « les extremités » ; la
leçon du feuillet est gardée. De même, « la seconde par
$e^{(1-x)\omega}$ » est écrit ainsi, alors que la réduction qui précède
multiplie par $e^{(x-1)\omega}$. Un mot court biffé après le premier
« au lieu de » n'est pas lu. Les deux grandes expressions de
$\frac{\alpha}{\alpha'}$ sont écrites sur le feuillet comme un produit
de deux fractions écrites côte à côte, séparées d'un point ; elles sont
données ici avec un point de multiplication. Le verso (vue 298) est
blanc.}

\page{299}

23 A l'egard du cas ou l'une et l'autre extrémités sont appuiées
Euler trouve $\frac{\alpha}{\alpha'} = \frac{e^{x\omega} - e^{(2-x)\omega}}{e^{x\omega} - e^{-x\omega}}$ et il est clair que l'équation
ne mène pas a conclure $\alpha = \alpha'$, car rien n'engage a \struck{conclure} faire
$e^{x\omega} - e^{(2-x)\omega} = e^{x\omega} - e^{-x\omega}$, mais il faut remarquer qu'ici tout
les coëfficiens a l'exception de $\delta$, tout \struck{en vertu} nuls, en
vertu de l'état des extremités, de sorte que $\frac{\alpha}{\alpha'}$ est essentielle-
ment une quantité indéterminée et que parconséquent on
ne peut rien conclure de l'équation $\frac{\alpha}{\alpha'} = \frac{e^{x\omega} - e^{(2-x)\omega}}{e^{x\omega} - e^{-x\omega}}$,

On voit en effet qu'après avoir trouvé pour la première solution
qui comme je l'ai déjà remarqué est la seule qui s'accorde avec
l'état des extremités, et dans la supposition $x = \frac{1}{2}$
\[ \alpha = e^{\frac{1}{2}\omega} - e^{\frac{3}{2}\omega};\quad \beta = e^{\frac{3}{2}\omega} - e^{\frac{1}{2}\omega};\quad \gamma = -\frac{(e^{\frac{1}{2}\omega} - e^{\frac{3}{2}\omega})(e^{\frac{1}{2}\omega} - e^{-\frac{1}{2}\omega})}{\text{sin}\,\frac{1}{2}\omega};\quad \delta = 0. \]
\[ \alpha' = e^{\frac{1}{2}\omega} - e^{-\frac{1}{2}\omega},\quad \beta' = -e^{2\omega}(e^{\frac{1}{2}\omega} - e^{-\frac{1}{2}\omega});\quad \gamma' = -\frac{(e^{\frac{1}{2}\omega} - e^{-\frac{1}{2}\omega})(e^{\frac{1}{2}\omega} - e^{\frac{3}{2}\omega})}{\text{sin}\,\frac{1}{2}\omega - \text{cos}\,\frac{1}{2}\omega\,\text{tang.}\,\omega};\quad \delta' = \frac{(e^{\frac{1}{2}\omega} - e^{-\frac{1}{2}\omega})(e^{\frac{1}{2}\omega} - e^{\frac{3}{2}\omega})\,\text{tang.}\,\omega}{\text{sin}\,\frac{1}{2}\omega - \text{cos}\,\frac{1}{2}\omega\,\text{tang.}\,\omega} \]
il multiplie ensuite tous ces coëfficiens par $\text{sin}\,\frac{1}{2}\omega$ afin de les réduire
a leur véritable valeur, c'est a dire, a celle qui convient a l'état
dela question; ce qui lui donne:
\[ \alpha = \alpha' = 0,\quad \beta = \beta' = 0,\quad \gamma' = -\gamma \quad\text{et}\quad \delta = \delta' = 0, \]
conformement a ce que j'ai \struck{trouvé} devoir avoir lieu toutes
les fois que l'on aiégard a l'état des extremités, on a seulement
$\delta = -\delta'$ au lieu de $\delta = \delta'$ parcequ'on atribue le changement
de signe au coëfficient $\delta$ au lieu de le considerer comme
appartenant a l'ordonnée $y$ qui est prise de l'autre côté de
l'axe.

On voit donc que tout ce qu'on peut savoir sur l'effet d'un
point d'appui placé dans un \struck{lieu} point quelconque de l'étendue
dela lame \struck{est} resulte dela supposition $y = 0$ dans ce
point combiné avec les conditions dependantes de l'état
des extrèmités et que l'appareille des quatre équations
trouvées par Euler est au moins inutile puisqu'après delong
\struck{de} calculs ont atteint aucun resultat nouveau \struck{et que deplus}
ce qui doit necessairement avoir lieu \uncertain{par} cause que les deux
\note{En tête, au milieu, « 38 » souligné, le chiffre repassé et empâté ;
en haut à droite « 146 », à l'encre. Le paragraphe porte le numéro 23.
Les huit coefficients cités sont, au signe et à la disposition près, ceux
du § 50 du mémoire d'Euler, qu'elle a copiés en latin à la vue 200 de ce
même volume. « aiégard » est écrit en un mot, pour « a égard » ;
« l'appareille » pour « l'appareil » ; « tout » pour « tous ». Le $\gamma$
est tracé comme un « v ». Le feuillet s'arrête sur « les deux » ; la
suite n'est pas dans ce lot. Le verso (vue 300) est blanc.}

\end{document}
